Reported rates of ADHD \parencite{cuiDSM5ChangesCOVID192026,mckechnieAttentiondeficitHyperactivityDisorder2023}  and autism \parencite{zeidanGlobalPrevalenceAutism2022} have risen sharply. Concurrently, mental health awareness campaigns have proliferated \parencite{bolinskiEffectEmentalHealth2020}, and mental health-related discourse has become increasingly common in everyday life \parencite{foulkesAreMentalHealth2023}, giving rise to what has been described as \textit{therapy-speak}. Therapy-speak refers to the imprecise and superficial integration of psychotherapy language into everyday communication \parencite{isern-masUnmaskingTherapyspeak2025}.

This concurrence has led some scholars to question the temporal relationship between rising prevalence rates and the increasing everyday use of mental health language. \textcite{foulkesAreMentalHealth2023} describe this as the \textit{Prevalence Inflation Hypothesis}, which proposes a cyclical, intensifying relationship between the two phenomena. On one side of this cycle, rising prevalence rates may understandably increase mental health-related discourse in everyday life, hereafter referred to as therapy-speak. On the other side, the spread of therapy-speak may itself contribute to further increases in prevalence rates through overdiagnosis and overinterpretation \parencite{foulkesAreMentalHealth2023}.

Before considering the potential risks of therapy-speak, its positive effects should be emphasised. The popularisation of therapy-speak may give people language for subtle emotional states \parencite{page18HowHarnessPower}, improve mental health literacy \parencite{schomerusEvolutionPublicAttitudes2012}, and thereby reduce stigma around mental illness \parencite{sampognaImpactSocialMarketing2017, flearyRelationshipHealthLiteracy2022}. It may also support better recognition and more accurate reporting of mental health problems \parencite{foulkesAreMentalHealth2023} and facilitate online identity and community formation, particularly around ADHD and autism \parencite{ginappExperiencesAdultsADHD2023}.

At the same time, therapy-speak carries several risks that are predominantly epistemic in nature. These risks are not mutually exclusive and may overlap. First, therapy-speak may encourage overinterpretation, whereby milder and more transient forms of distress are conflated with mental health problems, often in conjunction with self-diagnosis \parencite{foulkesAreMentalHealth2023}. At best, this may alter self-concept and behaviour \parencite{foulkesAreMentalHealth2023,alperTikTokAlgorithmicallyMediated2025}. At worst, when an individual interprets and labels their psychological experiences as a mental health problem, this may bring those symptoms into existence in the manner of a self-fulfilling prophecy. Adolescents may be particularly vulnerable to these dynamics because they are more prone to rumination, peer influence, and media messaging \parencite{foulkesAreMentalHealth2023}.

Second, therapy-speak may contribute to the psychiatrisation of everyday suffering and distress (Brinkmann, 2014; Beeker et al., 2024). Third, mental health problems may be glamorised or romanticised, particularly among adolescents on social media \parencite{ndourRomanticisationMentalHealth2025}. Fourth, therapy-speak may spread inaccurate psychological information, including inaccurate information about autism \parencite{aragon-guevaraReachAccuracyInformation2025} and ADHD \parencite{devriesExploringConceptCreep2025}. Finally, therapy-speak may erode the meaning and relevance of mental-health-related terms, a process that can also be understood as trivialisation. This can contribute to hermeneutical injustice by depriving people who actually live with a certain mental condition of the words they need to describe their experiences \parencite{isern-masUnmaskingTherapyspeak2025} and by reducing their symptoms to mere personality traits, thereby denying them a fully recognised psychiatric identity \parencite{spencerIsntEveryoneLittle2021}.

This latter concern---the erosion of meaning in psychotherapy terms---is the focus of the present study, which investigates diachronic lexical semantic change in the terms ADHD and autism. These two terms stand out, arguably more than any others, when considering psychotherapy-related terms that have permeated mainstream discourse \parencite{page18HowHarnessPower}. From January to May 2024, ADHD was the subject of 25,080 media articles, compared with 5775 articles during the equivalent period in 2014 \parencite{martinChangingPrevalenceADHD2025}. In May 2026, adhd and autism had been hashtagged in more than 5.2 million and 3.9 million videos on TikTok, respectively.

The following literature review first synthesises empirical evidence on lexical semantic change (LSC) in mental-health-related terms over recent decades, before reviewing the computational approaches used to study these processes and identifying the gap addressed by this project.
